% =====================================================================
% Section 7 -- Conclusions (file name kept from the earlier numbering).
% Numbers are macros only.
% =====================================================================
\section{Conclusions}\label{sec:conclusions}

The longitudinal dynamics of a belt conveyor loop with two wave speeds, a gravity take-up
modelled as a pulley with mass anywhere on the return strand and a drive anywhere on the loop
reduce to a closed-form characteristic equation. The free end that a light take-up is known to
form holds at any position, and the take-up position sets the lengths of the two strands into
which it splits the loop.

The take-up mass enters as a tip mass of a quarter of its value. Because it follows from the
take-up tension, it lengthens the fundamental period of the published conveyors of
\qty{1}{\kilo\metre} or longer by at most \qty{\NumLongShiftMax}{\percent}, and
$\Tt\lesssim\Wret/4$ suffices to neglect it in that period for a counterweight hung directly on
two strands near the drive of a loaded belt ($\gamma\ge1.4$).

Moving the take-up from the head to the tail shortens the fundamental period by a factor
\NumHeadTailRatioGammaThree--\NumHeadTailRatioGammaOneFive\ on loaded belts ($\gamma=1.5$--3)
and up to
\NumHeadTailRatioGammaOne\ on an empty one, so relocating it in a model is not an equivalent
transformation. Each strand starts up along a universal curve, a finite d'Alembert sum, which
gives the drive tensions, the take-up travel and the grip condition.

On three published conveyors, with hypothetical speed-controlled starts, the take-up position
acts mainly through the statics, and dynamically through the start time. With a single drive
pulley the grip fixes the shortest start: the \qty{30}{\second} starts of \citet{lodewijks1996}
would have slipped, as a rigid-body balance already shows. A start record with a free take-up,
a speed-controlled drive and published conveyor data would test the model directly.
